\documentclass[11pt]{article}
\usepackage[margin=1in]{geometry}
\usepackage[T1]{fontenc}
\usepackage[utf8]{inputenc}
\usepackage{parskip}
\usepackage{hyperref}

\begin{document}
\noindent
\textbf{In-Gee Kim}\\
Bona Sapiens, Inc.\\
Seoul 07332, Republic of Korea\\
\href{mailto:kim.ingee@bonasapiens.com}{kim.ingee@bonasapiens.com}

\vspace{1em}
\noindent\today

\vspace{1em}
\noindent
The Editors\\
\textit{Physical Review E}

\vspace{1em}
\noindent
\textbf{Re: Submission of manuscript ``Linear response representation of two-fermion plasmas''}

\vspace{0.5em}
Dear Editor,

We submit the accompanying manuscript for consideration as a Regular Article in \textit{Physical Review E}.
This paper does not merely present a new numerical scheme for the finite-temperature Lindhard linear response and two-component random-phase approximation (RPA) dielectric functions.
Rather, it exposes and strictly resolves a fundamental analytic breakdown---present in the literature for nearly 70 years---endemic to standard contour integration over Matsubara frequencies.

\textbf{Context and the 70-year inconsistency.}
For decades, the evaluation of $\varepsilon(q,\omega)$ under partial degeneracy has relied on hybrid mathematical routes.
We mathematically demonstrate that these widely accepted formulations---which evaluate the imaginary part in a grand-canonical prescription while forcing the real part through canonical contour sums or mismatched analytic continuations---inherently violate Jordan's Lemma and break exact causality.
This ensemble inconsistency produces persistent, unremovable $\mathcal{O}(1)$ errors in $\Re\chi^L$ that become catastrophically evident in the highly degenerate $T \to 0$ limit.

\textbf{Main results: the causality-constrained framework.}
To restore exact causality, this manuscript formally abandons the divergent Matsubara pole summations.
Instead, we map the calculation into a strictly constrained physical space: $\Im\chi^L$ is derived in closed form, and $\Re\chi^L$ is obtained exclusively as its Kramers--Kr\"onig Hilbert transform.
Implemented via a modified sinc-quadrature rule, our formulation possesses an exact theoretical truncation error bound of $\mathcal{O}(\sqrt{N} \exp(-\sqrt{\pi d \alpha N}))$, suppressing numerical noise to $\sim 10^{-259}$ (well below machine precision).

The robustness of this deterministic pipeline is validated through:
\begin{itemize}
\item exact derivative-free extraction of the plasmon dispersion and Landau damping;
\item automated satisfaction of four macroscopic dynamic structure factor sum rules to double-precision limits, without ad hoc regularizations;
\item seamless extension to the full two-component RPA susceptibility tensor ($S_{ee}$, $S_{ii}$, $S_{ei}$) capturing exact interspecies entanglement.
\end{itemize}

\textbf{Appropriateness for \textit{Physical Review E}.}
By identifying the topological and thermodynamic failures of legacy contour routes and providing a mathematically rigorous, deterministic solution, this work touches the very core of many-body theory, plasma physics, and spectral representation.
We anticipate this causality-constrained framework will establish a new fundamental baseline for evaluating dynamic structure factors in Fermi systems.

\textbf{Submission history and open-source verifiability.}
This manuscript has not been published previously and is not under consideration elsewhere.
To ensure absolute transparency and facilitate immediate referee verification of our $\sim 10^{-259}$ error bounds and exact sum-rule matching, the complete \texttt{C++20} MosaiQ\-Lindhard engine is open-sourced and available on GitHub/Zenodo (DOI pending).

\textbf{PhySH concepts (suggested).}
Plasma physics (primary); Multicomponent plasmas; Electron gas; Dielectric function; Dynamic structure factor; Random phase approximation; Linear response; Kramers--Kr\"onig relations.

We look forward to a rigorous review process.
Please contact me at \href{mailto:kim.ingee@bonasapiens.com}{kim.ingee@bonasapiens.com} regarding this submission.

\vspace{1em}
\noindent
Sincerely,

\vspace{1.5em}
\noindent
In-Gee Kim

\end{document}
